\chapter{Arithmetic realization and closure}\label{ch:arith}

We now study which sets of orderings SELECT-II can produce. Throughout, a \emph{content} is a set $C=\{c_0<\dots<c_{k-1}\}$ of distinct positive integers; a label is addressed by its \emph{rank} $r\in\{0,\dots,k-1\}$, an arrangement is a permutation of the ranks, and $A_C(m)\subseteq S_k$ is the set of arrangements admissible from the source $m$ (\cref{sec:complex}).

\section{The feasibility lemma}
For a set $I$ of ranks put $R(I)=\prod_{r\in I}\frac{(c_r+1)^2}{c_r}\in\Q$ (with $R(\emptyset)=1$) and let $D(I)$ be the denominator of $R(I)$ in lowest terms. For an arrangement $w=(w_1,\dots,w_k)$ with prefix sets $I_j=\{w_1,\dots,w_j\}$ define its \emph{threshold} $K(w)=\operatorname{lcm}_{1\le j\le k}D(I_j)$.

\begin{lemma}[feasible-family law]\label{lem:F}
Let $x$ be an integer and $N$ a label. Then $N\mid x$ iff $x\,R(N)\in\Z$. Consequently $w$ is admissible from $m$ iff $m\,R(I_j)\in\Z$ for every $j=1,\dots,k$, iff $D(I_j)\mid m$ for every $j$, i.e.
\[
A_C(m)=\{w:\ K(w)\mid m\}=\{w:\ I_j(w)\in\mathfrak F_C(m)\ \forall j\},\qquad \mathfrak F_C(m)=\{I:\ D(I)\mid m\}.
\]
\status{\Analytic; additionally checked by the manuscript verifier by simulating SELECT-II against the law on all $720$ arrangements for six sources and on all sources $1..199$ for the content $(2,3,6,7)$ (\texttt{ch4\_lemmaF})}{GIG-I v003, Lemma F}
\end{lemma}
\begin{proof}
Since $\gcd(N,(N{+}1)^2)=1$, $xR(N)=x(N{+}1)^2/N$ is an integer iff $N\mid x$. The state after $j$ steps, if all were applicable, is $mR(I_j)$, which depends on $I_j$ only. By induction on $j$, step $j{+}1$ is applicable from the integer $mR(I_j)$ iff $mR(I_{j+1})\in\Z$. Finally, writing $R(I)=a/D$ in lowest terms, $mR(I)\in\Z$ iff $D\mid ma$ iff $D\mid m$.
\end{proof}

Hence arithmetic admissibility is \emph{prefix-defined} in the sense of \cref{def:prefix-defined} with the family $\mathfrak F_C(m)$ (the last prefix set is $[k]$ itself; if $D([k])\nmid m$ the set $A_C(m)$ is empty, and for nonempty $A_C(m)$ the condition at $[k]$ is automatic). Everything in \cref{ch:swaps,ch:topology,ch:duality} that holds for prefix-defined sets therefore applies to SELECT-II.

\begin{proposition}[valuation form]\label{prop:valuation}
For a prime $p$ put $\omega_p(r)=v_p(c_r)-2v_p(c_r+1)$ and $\omega_p(I)=\sum_{r\in I}\omega_p(r)$. Then $v_p(D(I))=\max(0,-v_p(R(I)))=\max(0,\,-\!\sum_{r\in I}(2v_p(c_r{+}1)-v_p(c_r)))$ and the sign convention of the weights gives
\begin{gather*}
D(I)\mid m\iff \omega_p(I)\le v_p(m)\ \text{ for every prime }p,\\
\text{so}\qquad \mathfrak F_C(m)=\bigcap_p\{I:\ \omega_p(I)\le v_p(m)\}.\end{gather*}
\status{\Analytic}{GIG-I v005, PS-6}
\end{proposition}
\begin{proof}
$v_p$ is additive on products, $v_p\big(\frac{(c+1)^2}{c}\big)=2v_p(c{+}1)-v_p(c)=-\omega_p$, so $v_p(R(I))=-\omega_p(I)$ and the exponent of $p$ in the reduced denominator is $\max(0,\omega_p(I))$. Then $v_p(D(I))\le v_p(m)$ iff $\max(0,\omega_p(I))\le v_p(m)$ iff $\omega_p(I)\le v_p(m)$, because $v_p(m)\ge0$.
\end{proof}
A label contributes a positive weight at $p$ when $p\mid c_r$ (it \emph{consumes} $p$-divisibility of the state) and a negative weight $-2v_p(c_r{+}1)$ when $p\mid c_r+1$ (it \emph{supplies} $p$-divisibility). Only primes dividing some $c_r$ or $c_r+1$ matter.

\begin{corollary}[finite lattice of systems]\label{prop:finite-lattice}
Let $\mathrm{cap}_p=\max(0,\max_I\omega_p(I))$. Then $A_C(m)$ depends on $m$ only through $\big(\min(v_p(m),\mathrm{cap}_p)\big)_p$ over the finitely many relevant primes; hence the number of distinct sets $A_C(m)$, $m\in\N$, is at most $\prod_p(\mathrm{cap}_p+1)$.
\status{\Analytic}{GIG-INFO v001, GI-0}
\end{corollary}
\begin{proof}
If $v_p(m)\ge\mathrm{cap}_p$ the condition $\omega_p(I)\le v_p(m)$ holds for all $I$, so raising $v_p(m)$ beyond the cap changes nothing, and primes without weights impose nothing.
\end{proof}
For the content $\lambda_A=(2,3,4,5,8,9)$ the caps are $2^6,3^3,5^1$, giving at most $56$ systems (exactly $56$); for $\lambda_B=(2,3,6,7,8,9)$ the caps are $2^5,3^4,7^1$ (and $0$ at $5$), giving $60$. These are the two \emph{anchor contents} of \cref{ch:topology}. The source line collapses to a finite lattice of regimes: for example, ``the $60$ sources $288t$'' realise only $8$ distinct regimes for either anchor.

\section{Closure and minimal sources}
Define for $S\subseteq S_k$ the \emph{lcm} $L(S)=\operatorname{lcm}_{w\in S}K(w)$, with $L(\emptyset)=1$, and the \emph{closure} $\cl_K(S)=\{v\in S_k:\ K(v)\mid L(S)\}$. Call $S$ \emph{realizable} if $S=A_C(m)$ for some $m$. The family $\mathcal A(C)=\{A_C(m):m\in\N\}$ \emph{includes} $\emptyset$ when some $A_C(m)$ is empty.

\begin{theorem}[threshold closure]\label{thm:closure}
(i) $S$ is realizable iff $\cl_K(S)=S$. (ii) If $S$ is realizable, $L(S)$ divides every realizing source and itself realizes $S$: it is the unique divisibility-minimal (hence numerically smallest) realizing source. (iii) $\cl_K$ is a closure operator (extensive, monotone, idempotent) and $\cl_K(S)$ is the smallest realizable superset of $S$. (iv) $A_C(m)\cap A_C(m')=A_C(\gcd(m,m'))$, so $\mathcal A(C)$ is closed under intersection; the subfamily of \emph{nonempty} realizable sets is not. (v) A realizable $S$ is realized by a squarefree source iff $L(S)$ is squarefree.
\status{\Analytic; with the correction recorded in GIG-I v003 for the empty set (the nonempty-only family is not intersection-closed)}{GIG-I v002, T12; GIG-I v003 \texttt{CORRECTIONS/V002\_T12\_SCOPE}}
\end{theorem}
\begin{proof}
(i) If $S=A_C(m)$ then $K(w)\mid m$ for all $w\in S$, so $L(S)\mid m$; if $K(v)\mid L(S)$ then $K(v)\mid m$, so $v\in S$. Thus $\cl_K(S)\subseteq S$, and $S\subseteq\cl_K(S)$ trivially. Conversely if $\cl_K(S)=S$ then $A_C(L(S))=\{v:K(v)\mid L(S)\}=\cl_K(S)=S$.
(ii) The first half of (i) shows $L(S)\mid m$ for every realizing $m$; the second shows $L(S)$ realizes $S$.
(iii) Extensive: shown. Monotone: $S\subseteq S'\Rightarrow L(S)\mid L(S')\Rightarrow\cl_K(S)\subseteq\cl_K(S')$. Idempotent: every $K(v)$ with $v\in\cl_K(S)$ divides $L(S)$, so $L(\cl_KS)\mid L(S)$, and $S\subseteq\cl_KS$ gives $L(S)\mid L(\cl_KS)$; hence $L(\cl_KS)=L(S)$ and $\cl_K\cl_KS=\cl_KS$. If $T\supseteq S$ is realizable then $\cl_K(S)\subseteq\cl_K(T)=T$.
(iv) $w\in A(m)\cap A(m')$ iff $K(w)\mid m$ and $K(w)\mid m'$ iff $K(w)\mid\gcd(m,m')$.
(v) If a squarefree $m$ realizes $S$ then $L(S)\mid m$ is squarefree; conversely $m=L(S)$ then realizes $S$.
\end{proof}
\begin{example}[the erratum]\label{ex:erratum}
For $C=(2,3)$: $K((2,3))=\operatorname{lcm}(D(\{2\}),D(\{2,3\}))=\operatorname{lcm}(2,1)=2$ and $K((3,2))=3$. Hence $A(2)=\{(2,3)\}$, $A(3)=\{(3,2)\}$, $A(1)=\emptyset$, $A(6)=$ both, and $A(2)\cap A(3)=A(\gcd(2,3))=A(1)=\emptyset$: the intersection of two nonempty realizable sets is empty. This is why the family used for regimes is $\mathcal A(C)$ including $\emptyset$, with $L(\emptyset)=1$.
\end{example}

\section{Block factorization and the cube-product property}
\begin{theorem}[ordered-block factorization]\label{thm:block}
Let $A$ be prefix-defined by $\mathfrak F$ and let $F=(B_1|\dots|B_r)$ be a face of $P_k$. A vertex $w=\pi_1\cdots\pi_r$ of $F$ ($\pi_i$ an ordering of $B_i$) lies in $A$ iff for every $i$ every nonempty prefix set $U$ of $\pi_i$ satisfies $B_{<i}\cup U\in\mathfrak F$ (proper sets only). Hence
\begin{gather*}
V(F)\cap A=E_1\times\dots\times E_r\quad(\text{concatenation}),\\ E_i=\{\pi_i\in S_{B_i}:\ B_{<i}\cup U\in\mathfrak F\ \forall\ \emptyset\ne U\ \text{prefix of}\ \pi_i\},\end{gather*}
and $F\in K_A$ iff every $E_i$ is all of $S_{B_i}$, iff every Boolean interval $[B_{<i},B_{\le i}]$ lies in $\mathfrak F$. In particular $|V(F)\cap A|=\prod_i|E_i|$.
\status{\Analytic (for all $k$, all prefix-defined systems; for SELECT-II this is GIG-I v003 T16, and $E_i$ is determined by the formal entry state $s_i=m\,R(B_{<i})\in\Q$)}{GIG-I v003, T16; the generalisation to arbitrary prefix-defined $A$ is the observation of this monograph, with the same proof}
\end{theorem}
\begin{proof}
The proper prefix sets of $w$ are the sets $B_{<i}\cup U$ with $U$ a nonempty prefix set of $\pi_i$ (the set $B_{<i}\cup\emptyset=B_{\le i-1}$ is already the last prefix set of $\pi_{i-1}$). Membership of such a set in $\mathfrak F$ is a condition on $\pi_i$ alone, because $B_{<i}$ is the \emph{set} of earlier blocks and does not depend on the order chosen inside them. So $w\in A$ iff $\pi_i\in E_i$ for all $i$. If all $E_i$ are full then $V(F)\subseteq A$ iff all sets $B_{<i}\cup U$, $\emptyset\neq U\subseteq B_i$, lie in $\mathfrak F$ (every such $U$ is a prefix set of some ordering of $B_i$), i.e. each interval $[B_{<i},B_{\le i}]$ is contained in $\mathfrak F$.
\end{proof}

\begin{corollary}[cube-product property]\label{cor:CP}
Let $A$ be prefix-defined. (a) For a commuting cube $Q_3$ (face of type $(2,2,2)$) the set $V(Q_3)\cap A$ is a product $E_1\times E_2\times E_3$ with $E_i\subseteq S_2$, hence of size $0,1,2,4$ or $8$ and a coordinate sub-box; in particular a set of six vertices of a cube is never of this form. (b) \emph{Square closure}: if three vertices of a commuting square lie in $A$, so does the fourth. (c) Consequently $A$ has the \emph{cube-product property} (CP): if six vertices of a commuting cube lie in $A$, all eight do. (d) A Petrie $6$-cycle and a $2\times3$ ``domino'' of a cube are never the trace $V(Q_3)\cap A$.
\status{\Analytic; in the manuscript verifier $6{,}000$ random cube slices of random prefix families gave no violation (\texttt{ch5\_cube\_product\_property\_prefix\_defined})}{GIG-I v002, T11 and corollaries}
\end{corollary}
\begin{proof}
A product of three subsets of a two-element set has $|E_1||E_2||E_3|\in\{0,1,2,4,8\}$ elements; a product of two $2$-element sets and a $1$-element set is a coordinate rectangle of size $4$ and a subset of a square containing three corners has all four corners.
\end{proof}
\begin{remark}[repeated labels]\label{rem:T11}
The cubical factorization of GIG-I v002 (T11) is stated for words in which labels outside the exchanged pairs may repeat. The same argument applies: the \emph{formal} states $s_1=m_0,\ s_{i+1}=s_i\prod_{x\in u_i}R(x)\in\Q$ do not depend on orientation choices (multiplication in $\Q$ commutes), a single-letter block is locally admissible iff $s_i\in\Z$ and $x\mid s_i$, and a pair block admits an orientation iff the order is admissible step by step from $s_i$; by induction on the blocks the admissible orientations form a product $G\times E_1\times\dots\times E_d$, $G\in\{0,1\}$, $E_i\subseteq\{0,1\}$. It is used only for the statement about contents with repeated labels and is not needed below.
\end{remark}

\section{Face thresholds and the divisibility filtration}
\begin{theorem}[face thresholds, filtration]\label{thm:face-threshold}\label{thm:filtration}
Let $X_C(m)=K_{A_C(m)}$ be the history complex. (a) A face $F$ of $P_k$ lies in $X_C(m)$ iff $K_F\mid m$, where $K_F=\operatorname{lcm}\{D(I):I\in\Pref(F)\}=\operatorname{lcm}_{w\in V(F)}K(w)$; and $F'\subseteq F\Rightarrow K_{F'}\mid K_F$. (b) If $m\mid m'$ then $A_C(m)\subseteq A_C(m')$ and $X_C(m)\subseteq X_C(m')$; moreover $F\in X_C(m)$ iff $v_q(K_F)\le v_q(m)$ for every prime $q$. (c) For every inclusion $X\subseteq X'$ of such complexes and every $j$,
\[
\Delta\beta_j=\Delta f_j-\Delta\rho_j-\Delta\rho_{j+1},
\]
where $f_j$ counts $j$-faces, $\rho_j=\operatorname{rank}\partial_j$ over $\F$ and $\Delta$ means $X'$ minus $X$. In particular $\Delta\beta_2=\Delta f_2-\Delta\rho_2-\Delta\rho_3$: a cover creating no $2$- or $3$-faces leaves $\beta_2$ unchanged.
\status{\Analytic}{GIG-I v003, T17, T18}
\end{theorem}
\begin{proof}
(a) By \cref{lem:faces-prefix} and \cref{lem:F}, $F\in X_C(m)$ iff $D(I)\mid m$ for all $I\in\Pref(F)$ iff $K_F\mid m$. The set of all prefix sets of vertices of $F$ is $\Pref(F)$, so $K_F=\operatorname{lcm}_wK(w)$ (up to the global factor $D([k])$, which divides every $K(w)$). If $F'$ is a face of $F$ then $V(F')\subseteq V(F)$, so $K_{F'}\mid K_F$. (b) $D(I)\mid m\mid m'$ gives $D(I)\mid m'$; the valuation statement is divisibility prime by prime. (c) Rank--nullity gives $\beta_j=f_j-\rho_j-\rho_{j+1}$ for each complex; subtract.
\end{proof}
The complexes $X_C(m)$ thus form a multi-parameter filtration indexed by the exponent vectors $(v_q(m))_q$: the \emph{source filtration at fixed content}, not the evolution of a state under SELECT-II. Events of this filtration (a face being born as $m$ becomes more divisible) are the topological ``activations'' studied in \cref{ch:topology}.

\section{Counting laws on the anchors}
Define the regime sizes $N(m)=|A_C(m)|$. For $\lambda_A$ at $m=288t$ the eight regimes with $A_C(288)\subseteq A_C(288t)$ have sizes $372,408,444,480,612,648,684,720$, and for $\lambda_B$ at $m=432t$ sizes $276,312,324,360,612,648,684,720$.
\begin{observation}[counting laws]\label{obs:counting}
For $\lambda_A$, $N(288t)=372+36[2\mid t]+72[3\mid t]+240[5\mid t]$ for all $t=1,\dots,60$. For $\lambda_B$, $N(432t)=276+36[2\mid t]+48[3\mid t]+336[7\mid t]$ fails additivity at $t=21$ and $t=42$ (the multiples of $21$ up to $60$): $N(432\cdot21)=684$ while the additive prediction is $660$, an interaction of $+24$.
\status{\FiniteExact; $t\le60$, re-derived by the manuscript verifier (\texttt{ch4\_count\_laws}); no proof of the general form is claimed}{GIG-INFO v001, GI-6}
\end{observation}
These laws are used in \cref{ch:separation}: on the comparable window the history count $N$ is injective in the regime, which is why $N$ alone recovers the hole state there.
